\section{Julia-Mengen und Fatou-Mengen}
\label{sec:julia}
Julia-Mengen und Fatou-Mengen entstehen aus der Überlegung was geschieht, wenn man eine Funktion $f$ wiederholt auf ein $z \in \mathbb{C}$ anwendet.
Hierraus ergibt sich für jedes $z$ eine Folge komplexer Zahlen
\begin{equation}
    z \rightarrow f(z) \rightarrow f(f(z)) \rightarrow \ldots \:,
    \label{julia_folge}
\end{equation}
bzw.
\begin{equation*}
    z_{n + 1} = f(z_n) \quad \text{mit}\: n \in \mathbb{N}_0\: \text{und einem Startwert}\: z_0 \in \overline{\mathbb{C}} = \mathbb{C} \cup \infty.
\end{equation*}

Je nach Wahl des Startwerts $z_0$ kann diese Folge zwei Verhalten zeigen:
\begin{enumerate}
    \item Eine kleine Änderung des Startwerts führt zur praktisch derselben Folge. Der Wert $z_0$ wird der Fatou-Menge $F_f$ von $f$ zugeordnet.
    \item Eine kleine Änderung des Startwerts führt zu einer komplett anderen Folge. In diesem Fall wird $z_0$ der Julia-Menge $J_f$ von $f$ zugeordnet.
\end{enumerate}

\subsection{Periodische Folgen}
Das Verhalten von solchen Folgen ist besonders dann interessant, wenn diese periodisch sind, bzw. $z_n = z_0$ für ein $n > 0$ gilt.
In dem Fall wird $z_0$ als \emph{periodischer Punkt} bezeichnet.
Ist $n$ die kleinste natürliche Zahl mit dieser Eigenschaft, so heisst $n$ die \emph{Periode} des Zyklus.
Falls dies bereits für $n = 1$ zutrifft, wenn also $f(z_0) = z_0$ gilt, so ist $z_0$, nicht nur ein periodischer Punkt, sondern auch ein \emph{Fixpunkt} von $f$.
Aus diesen Definitionen geht hervor, dass ein periodischer Punkt von $f$ mit Periode $n$ gleichzeitig ein Fixpunkt der $n$-fach zusammengesetzten Funktion $f^n$ ist.

Um die Komplexität der Fragestellung zu reduzieren, werden zunächst nur Fixpunkte betrachtet.

\subsubsection{Erstes Beispiel}
Am Beispiel der Funktion 
\begin{equation*}
    f(z) = z^2,
\end{equation*}
lassen sich solche periodische Folgen beobachten.
Durch $n$-maliges Anwenden erhält man $f^n(z) = z^{2^n}$.
Weiter hat die Funktion $f$ drei Fixpunkte: $0, 1$ und $\infty$, denn für diese Punkte gilt $f(z_0) = z_0$.

Entwickelt man nun die Folge \eqref{julia_folge}, ausgehend von den Startwerten $0$ und $\infty$, so erhält man $f^n(0) = 0$ bzw. $f^n(\infty) \to \infty$.
Nun werden diese Startwerte leicht verändert, dies geschieht mittels einem $\varepsilon \ll 1$:
\begin{itemize}
    \item Für $z_0 = 0$: $f^n(0 + \varepsilon) = (0 + \varepsilon)^{2^n} \to 0$ mit $|\varepsilon| > 0$
    \item Für $z_0 \to \infty$: $f^n(\infty + \varepsilon) = (\infty + \varepsilon)^{2^n} \to \infty$ mit $|\varepsilon| > 0$
\end{itemize}
Obwohl die Startwerte verändert wurden, erhält man schlussendlich die gleichen Resultate, wie wenn man die Startwerte nicht verändert hätte.
Anders ausgedrückt führen kleine Änderungen der Startwerte zu praktisch denselben Folgen.
Es scheint, dass diese beiden Startwerte Folgen in ihrer Umgebung anziehen.

Für den Startwert $1$ erhält man $f^n(1) = 1$, durch leichtes Verändern jedoch:
\begin{itemize}
    \item $f^n(1 + \varepsilon) = (1 + \varepsilon)^{2^n} \to \infty$ mit $\varepsilon > 0$
    \item $f^n(1 - \varepsilon) = (1 - \varepsilon)^{2^n} \to 0$ mit $\varepsilon > 0$
\end{itemize}
Hier scheint es also, dass eine kleine Änderung des Startwerts zu einer ganz anderen Folge führt, dieser Startwert stösst Folgen in seiner nahen Umgebung weg.

\subsubsection{Zweites Beispiel}
Als zweites Beispiel dient die Funktion
\begin{equation*}
    h(z) = z + z^7.
\end{equation*}
Diese Funktion hat Fixpunkte bei $-\infty, 0$ und $\infty$.
Die Fixpunkte $-\infty$ und $\infty$ verhalten sich erneut anziehend und werden daher nicht näher untersucht.
Nimmt man jedoch $0$ als Startwert, passiert etwas anderes.
Zunächst ergibt sich $h^n(0) = 0$, durch eine kleine Änderung erhält man,
\begin{equation*}
    h(0 + \varepsilon) = (0 + \varepsilon) + (0 + \varepsilon)^7 \text{ mit } 0 < |\varepsilon| \ll 1.
\end{equation*}
Da es sich bei $\varepsilon$ um eine sehr kleine Zahl handelt, verschwindet der zweite Summand und übrig bleibt der lineare Term.
Somit ist
\begin{equation*}
    h(0 + \varepsilon) = 0 + \varepsilon.
\end{equation*}
Wendet man nun erneut $h$ auf dieses Resultat an, erhält man wieder $0 + \varepsilon$.
Auch durch $n$-maliges Anwenden von $h$ ändert sich das Resultat nicht.
Die Folge bewegt sich also, anders als vorhin, nicht.
Der Startwert $0$, scheint also Folgen weder abzustossen, noch anzuziehen.

\subsubsection{Arten von periodischen Punkten}
\label{arten_periodische_punkte}
Das Verhalten einer Folge hängt also von der Art der periodischen Punkte ab.
Sei $z_0$ ein periodischer Punkt von $f$ mit Periode $n$, demzufolge ist $z_0$ ein Fixpunkt von $g = f^n$.
Dadurch wird die Frage nach dem Verhalten in der Nähe eines periodischen Punktes reduziert, auf die Frage nach dem Verhalten in der Nähe eines Fixpunktes.
Diese Frage lässt sich mithilfe der Taylorreihe um $z_0$
\begin{equation*}
    g(z) \approx z_0 + g^\prime(z_0) (z - z_0)
    \label{taylor}
\end{equation*}
beantworten, wobei höhere Potenzen in der Nähe von $z_0$ verschwinden.
Verwendet man $z_0$ als Startwert und ändert diesen leicht, bedeutet dies
\begin{equation*}
    z = z_0 + \varepsilon \text{ mit } |\varepsilon| > 0.
\end{equation*}
Setzt man dies in die Taylorreihe ein, erhält man
\begin{equation*}
    g(z) = g(z_0 + \varepsilon) \approx z_0 + g^\prime(z_0) \left((z_0 + \varepsilon) - z_0\right) = z_0 + g^\prime(z_0) \varepsilon.
\end{equation*}
Wird $g$ erneut und wiederholt auf dieses Resultat angewendet, so wird das Verhalten der entstehenden Folge einzig von $g^\prime(z_0)$ festgelegt.
Ist $|g^\prime(z_0)| < 1$, so wird der zweite Summand durch wiederholte Anwendungen immer kleiner und die Folge kehrt zu $z_0$ zurück.
Bei $|g^\prime(z_0)| = 0$ geschieht dies bereits nach der ersten Anwendung von $g$.
Im Fall dass $|g^\prime(z_0)| > 1$ ist, wird der zweite Summand immer grösser und die Folge kehrt nicht zu $z_0$ zurück, sondern konvergiert zu einem anderen Wert.

Wenn $|g^\prime(z_0)| = 1$ ist, ergibt sich ein Spezialfall.
Durch wiederholtes Anwenden scheint das Resultat immer $z_0 + \varepsilon$ zu bleiben, das stimmt aber nicht.
Die Formel \eqref{taylor} für die Taylorreihe ignoriert die höheren Potenzen, da diese in den drei anderen Fällen nichts am Ausgang ändern.
Wiederholtes Anwenden der Funktion führt dazu, dass die höheren Potenzen irgendwann nicht mehr verschwinden, und im Fall $|g^\prime(z_0)| = 1$ werden sie sogar sehr wichtig.
Das Verhalten der Folge wird hier durch diese höheren Potenzen bestimmt und kann, sowohl auf die eine, als auch auf die andere Art ausfallen.
Dieses Phänomen wird im Abschnitt \ref{sec:indifferent} näher beschrieben.

Somit sei nun
\begin{equation*}
    \lambda = g^\prime(z_0) = \left(f^n\right)^\prime(z_0).
\end{equation*}

Dann heisst der periodische Punkt $z_0$
\begin{itemize}
    \item \emph{stark anziehend}, wenn $\lambda = 0$,
    \item \emph{anziehend}, wenn $0 < \left|\lambda\right| < 1$,
    \item \emph{indifferent}, wenn $\left|\lambda\right| = 1$,
    \item \emph{abstosstend}, wenn  $\left|\lambda\right| > 1$.
\end{itemize}

Wendet man dies auf das vorherige Beispiel der Funktion $f(z) = z^2$ bzw. $f^n(z) = z^{2^n}$ an, bekommt man
\begin{equation*}
    \lambda = \left(f^n\right)^\prime(z_0) = 2^n z_0^{2^n - 1}.
\end{equation*}
Für die Fixpunkte von $f$ ergibt dies
\begin{itemize}
    \item $\lambda_0 = 0$,
    \item $\lambda_1 = 2^n$.
\end{itemize}
Somit sind diese Punkte \emph{stark anziehend}, bzw. \emph{abstossed}.

Im Fall des Fixpunkts $\infty$ ist zunächst ein Koordinatenwechsel nötig, da sich dieser --- ausserhalb der komplexen Zahlenebene --- auf der Riemannschen Zahlenkugel befindet.
Sei
\begin{equation*}
    \omega = \frac{1}{z},
\end{equation*}
damit ist
\begin{equation*}
    f(\omega) = \omega^2 = \frac{1}{z^2}.
\end{equation*}
Für $\omega \to 0$ verhält sich $f$ gleich, wie für $z \to \infty$.
Berechnet man nun $\lambda$, unter Berücksichtigung der neuen Koordinate, bekommt man
\begin{equation*}
    \lambda = \left(f^n\right)^\prime(\omega_0) = 2 \omega_0,
\end{equation*}
was für $\omega \to 0$, den Wert $0$ ergibt.
Insofern ist auch dieser Fixpunkt \emph{stark anziehend}.

Die Namensgebungen decken sich also mit dem beobachteten Verhalten von Folgen, ausgehend von diesen Startwerten.
Es scheint als ob eine Definition für die Julia-Mengen auf \emph{abstossenden} Fixpunkten aufbauen muss. 

\subsection{Beobachtungen}
\label{sec:beobachtungen}
Wie die vorhergehenden Beispiele zeigen, scheinen, per Definition aus Abschnitt \ref{sec:julia}, abstossende Fixpunkte zur Julia-Menge der jeweiligen Funktion zu gehören.
Ebenfalls scheinen anziehende Fixpunkte zur Fatou-Menge zu gehören.
Dies reicht jedoch noch nicht zur Definition dieser Mengen aus denn indifferente Punkte, nicht-fixe periodische Punkte und nicht-periodische Punkte wurden noch nicht genauer betrachtet.

\subsubsection{Indifferente Punkte}
\label{sec:indifferent}
Im Abschnitt \ref{arten_periodische_punkte} wurde beschrieben, dass indifferente periodische Punkte einen Spezialfall bilden.
Abhängig von den höheren Potenzen in der Taylorreihe, können indifferente Punkte entweder zur Julia-Menge oder zur Fatou-Menge ihrer jeweiligen Funktion gehören.
Wenn diese höheren Potenzen bei $n$-maliger Anwendung von $f$ verschwinden, gehört der indifferente Punkt zur Fatou-Menge.
Wachsen diese jedoch langsam, so verhält sich der indifferente Punkt wie ein abstossender Punkt und gehört daher zur Julia-Menge.
Weitere Details sind im Abschnitt 13 in \cite{julia:indifferente_punkte} beschrieben.

\subsubsection{Nicht-fixe periodische Punkte}
\label{sec:nicht_fix_periodisch}
Am Ende des Abschnitts \ref{arten_periodische_punkte} wurde festgestellt, dass die Funktion $f(z) = z^2$ einen abstossenden Fixpunkt hat.
Diese Funktion hat jedoch weitere abstossende, nicht-fixe periodische Punkte --- sogar unendlich viele.
Es handelt sich dabei um einige Einheitswurzeln, dies sind komplexe Zahlen, deren $m$-te Potenz $1$ ergibt.

Damit eine Einheitswurzel $z_e$, mit $z_e^m = 1$, periodisch ist, muss $f^n(z_e) = z_e^{2^n} = z_e$ für ein $n$ gelten.
Dividiert man die beiden rechten Terme der vorherigen Gleichung durch $z_e$ erhält man
\begin{equation*}
    z_e^{2^n - 1} = 1.
\end{equation*}
Nun muss $m\: |\: 2^n - 1$ bzw. $2^n - 1 \equiv 0 \mod(m)$ gelten.
Für alle ungeraden $m$ ist dies erfüllt, da $2^n \equiv 1 \mod(m)$.

Weiter sind Einheitswurzeln mit ungeraden Potenzen abstossend.
Dazu berechnet man
\begin{equation*}
    \lambda = \left(f^n\right)^\prime(z_e) = 2^n z_e^{2^n - 1}.
\end{equation*}
Weil $2^n - 1$ ein Vielfaches von $m$ ist, ist der zweite Faktor gleich $1$, es bleibt
\begin{equation*}
    \lambda = 2^n 1 = 2^n
\end{equation*}
und weil $\lambda > 1$, handelt es sich um \emph{abstossende} Punkte.

Einheitswurzeln liegen stets auf dem Einheitskreis, d.h. $|z_e| = 1$.
Wäre dies nicht der Fall, würde der Betrag durch das Potenzieren verschwinden, bzw. ansteigen und somit würde $z_e^m = 1$ nicht mehr gelten.
Daraus geht hervor, dass Einheitswurzeln mit ungerader Potenz auch zur Julia-Menge der Funktion $f$ gehören müssen.
Addiert man nämlich ein $\varepsilon$ mit $0 < |\varepsilon| << 1$ zu $z_e$, ist der Betrag entweder kleiner oder grösser als $1$.
Wiederholtes Anwenden von $f$ wird nun entweder gegen $0$ oder gegen $\infty$ gehen.
Ohne $\varepsilon$ bewirkt $f$, dass das Argument sich um den Einheitskreis dreht und periodisch zu $1$ zurückkehrt.

Somit sind nicht nur abstossende Fixpunkte, sondern auch abstossende periodische Punkte relevant zur Definition der Julia-Mengen.

\subsubsection{Nicht-periodische Punkte}
Punkte die nicht periodisch sind, können entweder zur Julia-Menge oder zur Fatou-Menge gehören.
Zum Beispiel gehört der Punkt $z_0 = e^{2\pi i\sqrt{2}}$ zur Julia-Menge der Funktion $f(z) = z^2 = r e^{2\pi i(2 \varphi)}$, obwohl er nicht periodisch ist.
Durch $n$-maliges Anwenden wird $f$ zu $f^n(z) = z^{2^n} = r^{2^n} e^{2\pi i(2^n \varphi)}$.
Setzt man den Punkt $z_0$ ein, d.h. setzt $\sqrt{2}$ für $\varphi$ ein, müsste
\begin{equation*}
    2^n \sqrt{2} = \sqrt{2}
\end{equation*}
für ein $n$ gelten, damit $z_0$ periodisch ist.
Dies ist aber nur für $n = 0$ der Fall, was bedeutet dass $z_0$ nicht periodisch ist.

Nimmt man $z_0$ als Startwert für eine Folge und verändert diesen leicht, d.h. verwendet $\sqrt{2} + \varepsilon$ für $\varphi$ im Exponenten, erhält man eine ganz andere Folge.
Nach $n$-maligem Anwenden erhält man $2^n \sqrt{2} + 2^n \varepsilon$ anstatt $2^n \sqrt{2}$, die Folge verhält sich chaotisch.
Dieser Punkt gehört daher zur Julia-Menge von $f$.

Verwendet man für die gleiche Funktion $0.5$ als Startwert und verändert diesen leicht, so erhält man $f^n(0.5 + \varepsilon) = (0.5 + \varepsilon)^{2^n}$, was gegen $0$ konvergiert.
Die Folge für den unveränderten Startwert konvergiert ebenfalls gegen $0$, was bedeutet, dass dieser Punkt zur Fatou-Menge von $f$ gehört.

Nicht-periodische Punkte können sich also auch unterschiedlich verhalten und gehören daher entweder zur Julia-Menge oder zur Fatou-Menge.

\subsection{Definition Julia-Mengen}
Wie im Abschnitt \ref{sec:beobachtungen} gezeigt, scheint es, dass eine Definition für Julia-Mengen:
\begin{itemize}
    \item Auf \emph{abstossenden} periodischen Punkten (nicht nur Fixpunkten) aufbauen muss,
    \item sowie \emph{indifferente} periodische Punkte und nicht-periodische Punkte berücksichtigen muss.
\end{itemize}

Betrachtet man erneut die Funktion $f(z) = z^2$, so ist bereits bekannt, dass alle \emph{abstossenden} periodischen und Fixpunkte in der Julia-Menge enthalten sind.
Alle diese Punkte liegen auf dem Einheitskreis und wie im Abschnitt \ref{sec:nicht_fix_periodisch} beschrieben, gibt es unendlich viele davon.

Diese \emph{abstossenden} Punkte sind jedoch nicht dicht, umgangssprachlich ausgedrückt gibt es Lücken auf dem Einheitskreis, welche nicht zu dieser Punktemenge gehören.
Allerdings ist die Distanz von jeder Lücke zu einem \emph{abstossenden} Punkt arbiträr klein.
Obwohl die ``Lückenpunkte'' selber nicht periodisch sind, verhalten sie sich gleich wie \emph{abstossende} periodische Punkte weil sie so nahe bei diesen liegen.

Entwickelt man die Folge \eqref{julia_folge} ausgehend von einem ``Lückenpunkt'' $z_l$ und entwickelt dann eine zweite Folge ausgehend von $z_l + \varepsilon$, erhält man zwei komplett andere Folgen.
Für den Startwert $z_l$ führt wiederholtes Anwenden von $f$ auch dazu, dass sich das Argument um den Einheitskreis dreht, jedoch nie mehr zu $1$ zurückkehrt.
Für $z_l + \varepsilon$ geht wiederholtes Anwenden gegen $0$, wenn $\varepsilon$ negativ ist und gegen $\infty$, wenn $\varepsilon$ positiv ist.

Was die \emph{indifferenten} periodischen Punkte angeht, so hat die Funktion $f$ keine davon.
Allgemein verhalten sich diese jedoch auch \emph{abstossend} wenn sie in einer solchen ``Lücke'' liegen.
Ihr genaues Verhalten ist aber wesentlich komplexer und zur vollständigen Beschreibung wird auf Abschnitt 13 in \cite{julia:indifferente_punkte} verwiesen.

Abschliessend bedeutet das, dass die Julia-Menge von $f$ dem Rand des Einheitskreises entsprechen muss
\begin{equation*}
    J_f = \left\{z \in \overline{\mathbb{C}}\:\mid\:|z| = 1\right\},
\end{equation*}
bzw. allgemeiner
\begin{equation*}
    J_f \coloneqq \text{Abschluss}\left\{z \in \overline{\mathbb{C}}\:\mid\:z\text{ ist ein abstossender periodischer Punkt von f}\right\}.
\end{equation*}
Dabei bezieht sich \emph{Abschluss} auf den topologischen Abschluss \cite{julia:abschluss}, dies bewirkt, dass die ``Lückenpunkte'' auch in $J_f$ enthalten sind.

\subsection{Definition Fatou-Mengen}
Die Definition der Fatou-Mengen lässt sich ebenfalls aus dem obigen Beispiel $f(z) = z^2$ herleiten.
Wie bereits beschrieben, gehören alle Startwerte auf dem Rand des Einheitskreises zur Julia-Menge.
Für alle anderen Startwerte, d.h. $|z| \neq 1$, erhält man, durch leichtes Verändern, mehr oder weniger die gleichen Folgen.
Folglich lässt sich die Fatou-Menge als
\begin{equation*}
    F_f = \left\{z \in \overline{\mathbb{C}}\:\mid\:|z| \neq 1\right\}
\end{equation*}
definieren oder allgemein formuliert
\begin{equation*}
    F_f \coloneqq \overline{\mathbb{C}} \setminus J_f.
\end{equation*}
